\documentclass[11pt]{article}
\usepackage[utf8]{inputenc}
\usepackage{amsmath,amssymb}
\usepackage[margin=1in]{geometry}
\usepackage{hyperref}
\emergencystretch=3em

\title{Signed reflections of an amplitude on symmetric boxes}
\author{Claude, with Daniel Murfet}
\date{2026-09-07}

\begin{document}
\maketitle
\sloppy

\begin{abstract}
Astra~\#18/\#19's rank-2 extension, first unit. For an amplitude $\eta$ on the symmetric box
$(-1,1]^d$ the coordinate reflections produce the signed symmetrisation
$\eta_{\rm sym}(u)=\sum_{\sigma\in\{\pm1\}^d}\big(\prod_i\sigma_i^{h_i}\big)\eta(\sigma u)$
(\texttt{symAmp}, indexed by $\sigma\colon\mathrm{Fin}\,d\to\mathrm{Bool}$), which is continuous when
$\eta$ is, evaluates at the origin to $\eta(0)\prod_i(1+(-1)^{h_i})$ (\texttt{symAmp\_zero}), and
reduces the symmetric-box dressed integral exactly to the unit-box one:
\[
\int_{(-1,1]^d}\eta(x)\,x^he^{-\beta Nx^{2k}}\,dx=\int_{(0,1]^d}\eta_{\rm sym}(u)\,u^he^{-\beta Nu^{2k}}\,du
\]
(\texttt{integral\_symBox\_eq\_symAmp}). Zero \texttt{sorry}/\texttt{axiom}.
\end{abstract}

\section{The things formalised}
{\small
\begin{verbatim}
def sgn (b : Bool) : ℝ := if b then -1 else 1
def reflect (σ : Fin d → Bool) (u : Fin d → ℝ) : Fin d → ℝ := fun i => sgn (σ i) * u i
def symAmp (h : Fin d → ℕ) (η) (u) : ℝ :=
  ∑ σ : Fin d → Bool, (∏ i, sgn (σ i) ^ h i) * η (reflect σ u)
theorem continuous_symAmp (hη : Continuous η) : Continuous (symAmp h η)
theorem symAmp_zero : symAmp h η 0 = η 0 * ∏ i, (1 + (-1 : ℝ) ^ h i)
theorem symAmp_cons (a) (w) : symAmp h η (Fin.cons a w)
    = symAmp (Fin.tail h) (fun w' => η (Fin.cons a w')) w
      + (-1) ^ h 0 * symAmp (Fin.tail h) (fun w' => η (Fin.cons (-a) w')) w
theorem integral_Ioc_symm_add (hg : IntegrableOn g (Ioc (-1) 1)) :
    ∫ a in Ioc (-1) 1, g a = ∫ a in Ioc 0 1, (g a + g (-a))
theorem integral_symBox_eq_symAmp (β) : ∀ d h k N η, Continuous η →
    ∫ x in symBox d, η x * ((∏ i, x i ^ h i) * exp (-(β * N * ∏ i, x i ^ (2 * k i))))
      = ∫ u in unitBox d, symAmp h η u * ((∏ i, u i ^ h i) * exp (-(β * N * ∏ i, u i ^ (2 * k i))))
\end{verbatim}
}

\section{Mathematics}
Induction on the dimension with the Bochner peel on both boxes (unit~188). The inner symmetric
integral at fixed first coordinate $a$ is the $d$-dimensional one with amplitude $\eta(a,\cdot)$ at
parameter $Na^{2k_0}$, hence by the induction hypothesis a unit-box integral of the symmetrised
section; the peeled coordinate is handled by $\int_{(-1,1]}g=\int_{(0,1]}(g(a)+g(-a))\,da$, and
$(\pm a)^{h_0}$ produces the sign $(-1)^{h_0}$ in front of the reflected section, matching the
splitting of the sign sum along \texttt{Fin.consEquiv}.

\section{Paper-facing meaning}
For interior normal coordinates the dressed integral is the unit-box integral of the signed
symmetrisation, so the whole continuous-amplitude theory applies with $\eta_{\rm sym}$ in place of
$\eta$; the corner value $\eta_{\rm sym}(0)=\eta(0)\prod(1+(-1)^{h_i})$ shows that an odd exponent kills
the equal-ratio corner coefficient (next unit), not the integral.

\section{Lean notes}
\begin{itemize}
\item Sum over sign vectors: \texttt{Finset.prod\_univ\_sum} with \texttt{Fintype.piFinset\_univ} turns
$\prod_i\sum_b$ into $\sum_\sigma\prod_i$; \texttt{(Fin.consEquiv \_).sum\_comp} then
\texttt{Fintype.sum\_prod\_type} and \texttt{Fintype.sum\_bool} split the first coordinate.
\item \texttt{IntervalIntegrable.iff\_comp\_neg} has an autoParam; apply it with named
\texttt{(f := g) (a := -1) (b := 0)} and \texttt{.mp}.
\item Using the closure box \texttt{piBox d (Ioc (-c) 1)} with $c\in\{0,1\}$ for both integrability
lemmas leaves \texttt{Ioc (-0) 1}; rewrite \texttt{neg\_zero} before \texttt{exact}.
\end{itemize}

\end{document}
